\documentclass[11pt]{article}
\usepackage[a4paper,margin=2.2cm]{geometry}
\usepackage{xcolor}
\usepackage{fontspec}
\usepackage{xeCJK}
\usepackage{enumitem}
\setmainfont{Times New Roman}
\setCJKmainfont{SimSun}
\setlist[itemize]{leftmargin=1.4em,itemsep=0.35em,topsep=0.35em}
\setlength{\parindent}{0pt}
\setlength{\parskip}{0.55em}
\pagestyle{plain}
\begin{document}
{\color[HTML]{1F5FD2}\LARGE\bfseries 特殊儿童语言发展训练方案\par}
{\color[HTML]{667085}\small 星轨原创教育资源：用于家庭与学校共同制定支持计划。\par}
\vspace{0.8em}
{\color[HTML]{111827}\large\bfseries 训练重点\par}
\begin{itemize}
\item 语言发展训练应同时关注理解、表达、轮流、共同注意和功能性沟通，而不只是增加词汇数量。
\item 目标是帮助儿童能请求、拒绝、评论、求助和分享兴趣，让语言真正服务于生活。
\end{itemize}
{\color[HTML]{111827}\large\bfseries 家庭情境\par}
\begin{itemize}
\item 亲子共读时先描述图片，再停顿等待儿童回应，并在儿童已有表达上做一点扩展。
\item 把喜欢的物品放在看得到但拿不到的位置，制造自然请求机会，同时接纳口语、手势、图片或 AAC。
\end{itemize}
{\color[HTML]{111827}\large\bfseries 学校情境\par}
\begin{itemize}
\item 上课前预教关键词，使用主题词板，并提供“我想要”“轮到我”“请帮忙”等句式起点。
\item 教师应记录儿童在哪些活动中最愿意表达，再把高动机活动安排进下一轮训练。
\end{itemize}
\vfill
{\color[HTML]{667085}\footnotesize 本材料为应用内原创教育内容，不能替代医疗、法律或正式评估建议。\par}
\end{document}
